% !TEX root = ../main.tex
\corrigesChapitre{15}{Mouvement dans un champ uniforme}
\corrigesSouples

\begin{blocLivre}{cor-15-1}
\begin{corrige}{1}{Une chute libre verticale}
\demarche{$\vec{a}=\vec{g}$, puis deux primitives avec $v_z(0)=0$ et $z(0)=0$.}
\begin{questions}
\item Dans le référentiel terrestre, galiléen, seul le poids agit : $m\,\vec{a}=m\,\vec{g}$, d'où $\res{a_z=+g}$ (axe vers
le bas).
\item $\res{v_z=g\,t}$ et $\res{z=\frac12\,g\,t^2}$.
\item $z=\qty{20}{m}$ pour $t=\sqrt{\dfrac{2\times20}{9{,}81}}=\res{\qty{2,0}{s}}$ ; alors $v=g\,t=\res{\qty{19,8}{m.s^{-1}}}$.
\item $m\,g\,h=\frac12\,m\,v^2$, d'où $v=\sqrt{2\,g\,h}=\qty{19,8}{m.s^{-1}}$ : même résultat.
\end{questions}
\end{corrige}
\end{blocLivre}
\begin{blocLivre}{cor-15-2}
\begin{corrige}{2}{Le champ d'un condensateur plan}
\demarche{$E=U/d$ ; la force électrique vaut $e\,E$.}
\begin{questions}
\item $\vec{E}$ perpendiculaire aux plaques, dirigé de la plaque positive vers la plaque négative,
$E=\dfrac{500}{0{,}020}=\res{\qty{2,5e4}{V.m^{-1}}}$.
\item $E$ \res{\text{double}} si $U$ double ; il est \res{\text{divisé par deux}} si $d$ double.
\item $F=e\,E=\res{\qty{4,0e-15}{N}}$ et $P=m\,g=\res{\qty{8,9e-30}{N}}$ : le poids est
\res{\text{totalement négligeable}} (rapport de l'ordre de $10^{14}$).
\end{questions}
\end{corrige}
\end{blocLivre}
\begin{blocLivre}{cor-15-3}
\begin{corrige}{3}{La bille qui quitte la table}
\demarche{Vitesse initiale horizontale ($\alpha=0$), hauteur initiale $h=\qty{0,80}{m}$.}
\begin{questions}
\item $x=2{,}5\,t$ et $z=-4{,}905\,t^2+0{,}80$ (en m et s).
\item $z=0$ pour $t=\sqrt{\dfrac{2\times0{,}80}{9{,}81}}=\res{\qty{0,40}{s}}$.
\item $x=2{,}5\times0{,}404=\res{\qty{1,0}{m}}$.
\item $v_z=-g\,t=\qty{-3,96}{m.s^{-1}}$ et $v=\sqrt{2{,}5^2+3{,}96^2}=\res{\qty{4,7}{m.s^{-1}}}$.
\end{questions}
\end{corrige}
\end{blocLivre}
\begin{blocLivre}{cor-15-4}
\begin{corrige}{4}{Les énergies d'une balle}
\demarche{$E_\mathrm{m}=E_\mathrm{c}+E_\mathrm{pp}$ se conserve : au sommet $v=0$, au sol $z=0$.}
\begin{questions}
\item $E_\mathrm{c}=\frac12\times0{,}20\times8{,}0^2=\res{\qty{6,4}{J}}$ ; $E_\mathrm{pp}=0{,}20\times9{,}81\times1{,}5=\res{\qty{2,9}{J}}$ ;
$E_\mathrm{m}=\res{\qty{9,3}{J}}$.
\item Seul le poids travaille (frottements négligés) : $E_\mathrm{m}$ \res{\text{se conserve}}.
\item Au sommet, $E_\mathrm{c}=0$ : $z_{\max}=\dfrac{E_\mathrm{m}}{m\,g}=\res{\qty{4,8}{m}}$.
\item Au sol, $E_\mathrm{pp}=0$ : $v=\sqrt{\dfrac{2\,E_\mathrm{m}}{m}}=\res{\qty{9,7}{m.s^{-1}}}$.
\end{questions}
\end{corrige}
\end{blocLivre}
\begin{blocLivre}{cor-15-5}
\begin{corrige}{5}{Accélérer un électron}
\demarche{Le travail de la force électrique vaut $e\,U$ ; théorème de l'énergie cinétique.}
\begin{questions}
\item $W=e\,U$ (positif : la force est dans le sens du mouvement).
\item $E_\mathrm{c}-0=e\,U$, d'où $E_\mathrm{c}=\num{1.60e-19}\times1000=\res{\qty{1,60e-16}{J}}$.
\item $v=\sqrt{\dfrac{2\,e\,U}{m}}=\res{\qty{1,9e7}{m.s^{-1}}}$.
\item $v$ est proportionnelle à $\sqrt{U}$ : il faut \res{\text{multiplier la tension par 4}} (\qty{4,0}{kV}).
\end{questions}
\end{corrige}
\end{blocLivre}
\begin{blocLivre}{cor-15-6}
\begin{corrige}{6}{Le lancer du poids}
\demarche{Équations du cours avec $v_0=\qty{13,5}{m.s^{-1}}$, $\alpha=42^\circ$ et $h=\qty{2,10}{m}$.}
\begin{questions}
\item $\vec{a}=\vec{g}$ est vertical : l'accélération et la vitesse initiale définissent un plan vertical, dont le centre
du poids ne sort pas.
\item $x=10{,}0\,t$, $z=-4{,}905\,t^2+9{,}03\,t+2{,}10$ ; $z=-0{,}0487\,x^2+0{,}900\,x+2{,}10$.
\item $z_{\max}=h+\dfrac{(v_0\sin\alpha)^2}{2g}=\res{\qty{6,3}{m}}$.
\item $z=0$ pour $t=\qty{2,05}{s}$ (racine positive), d'où une longueur de \res{\qty{20,6}{m}}.
\end{questions}
\end{corrige}
\end{blocLivre}
\begin{blocLivre}{cor-15-7}
\begin{corrige}{7}{Réussir son tir}
\demarche{Écrire que le point de coordonnées (4,20 ; 3,05) appartient à la trajectoire.}
\begin{questions}
\item $z=-\dfrac{g}{2\,v_0^2\cos^2\alpha}\,x^2+x\tan\alpha+h$.
\item En remplaçant $x$ et $z$ et en isolant $v_0^2$ : $\dfrac{g\,x^2}{2\,v_0^2\cos^2\alpha}=x\tan\alpha+h-z$, d'où la relation.
\item $v_0=\res{\qty{7,1}{m.s^{-1}}}$.
\item Avec $v_0=\qty{7,4}{m.s^{-1}}$, en $x=\qty{4,20}{m}$ : $z=\qty{3,41}{m}$ : le ballon passe \res{\text{au-dessus}} du
centre de l'anneau, de \qty{36}{cm}.
\end{questions}
\end{corrige}
\end{blocLivre}
\begin{blocLivre}{cor-15-8}
\begin{corrige}{8}{Dévier un faisceau d'électrons}
\demarche{$\vec{a}=-\dfrac{e}{m}\,\vec{E}$, constant : mouvement uniforme selon $x$, uniformément accéléré selon $y$.}
\begin{questions}
\item $E=\dfrac{U}{d}=\res{\qty{1,0e4}{V.m^{-1}}}$ ; la charge étant négative, la force est \res{\text{opposée à }\vec{E}},
dirigée vers la plaque positive.
\item $a=\dfrac{e\,E}{m}=\qty{1,76e15}{m.s^{-2}}$ ; $x=v_0\,t$ et $y=\frac12\,a\,t^2$, d'où
$y=\dfrac{e\,E}{2\,m\,v_0^2}\,x^2=2{,}20\,x^2$ (en mètres).
\item En $x=L$ : $y=2{,}20\times0{,}050^2=\res{\qty{5,5}{mm}}$, inférieur à $d/2=\qty{10}{mm}$ : les électrons
\res{\text{ne touchent pas}} les plaques.
\item $\tan\theta=\dfrac{v_y}{v_x}=\dfrac{a\,L}{v_0^2}=0{,}22$, soit $\res{\theta=12^\circ}$.
\end{questions}
\schema[corrige]{15-ex08}
\end{corrige}
\end{blocLivre}
\begin{blocLivre}{cor-15-9}
\begin{corrige}{9}{Tracer l'évolution des énergies}
\demarche{Vitesse par différences centrées, puis $E_\mathrm{c}=\frac12mv^2$ et $E_\mathrm{pp}=mgz$.}
\begin{questions}
\item En chaque point $M_i$, $v_x\approx\dfrac{x_{i+1}-x_{i-1}}{2\tau}$ et de même pour $v_z$ ; puis $v^2=v_x^2+v_z^2$.
\item L'origine est le sol ($z=0$). Les positions extrêmes n'ont pas de voisine des deux côtés : la vitesse ne peut pas y
être calculée par cette méthode.
\item $E_\mathrm{c}$ diminue (la balle monte et ralentit), $E_\mathrm{pp}$ augmente, et \res{E_\mathrm{m}\text{ reste constante}}
(\qty{1,29}{J}) : les frottements sont négligeables.
\item Stocker les valeurs dans trois listes (\code{lEc.append(Ec)}\,\dots), puis
\code{plt.plot(t, lEc)}, \code{plt.plot(t, lEp)}, \code{plt.plot(t, lEm)} et \code{plt.show()}.
\item Les frottements dissipent de l'énergie : $E_\mathrm{m}$ \res{\text{diminuerait}} au cours du temps.
\end{questions}
\schema[corrige]{rev-15-ex09-cor}
\legende{Points calculés avec les données du programme. Les écarts à la constance viennent des arrondis du pointage.}
\end{corrige}
\end{blocLivre}
\begin{blocLivre}{cor-15-10}
\begin{corrige}{10}{Le toboggan aquatique}
\demarche{Sans frottement, $E_\mathrm{m}$ se conserve ; avec frottement, $\Delta E_\mathrm{c}=W(\vec{P})+W(\vec{f})$.}
\begin{questions}
\item La réaction de la piste est perpendiculaire au mouvement et ne travaille pas : $v=\sqrt{2\,g\,h}=\res{\qty{12,5}{m.s^{-1}}}$.
\item $\qty{9,0}{m.s^{-1}}<\qty{12,5}{m.s^{-1}}$ : \res{\text{les frottements ne sont pas négligeables}}.
\item $\frac12\,m\,v^2-0=m\,g\,h+W(\vec{f})$, d'où $W(\vec{f})=1215-2354=\res{\qty{-1,1e3}{J}}$.
\item $W(\vec{f})=-f\,L$, d'où $f=\dfrac{1139}{30}=\res{\qty{38}{N}}$.
\end{questions}
\end{corrige}
\end{blocLivre}
\begin{blocLivre}{cor-15-11}
\begin{corrige}{11}{L'accélérateur linéaire}
\demarche{Chaque intervalle apporte $e\,U$ ; dans les tubes, le mouvement est uniforme.}
\begin{questions}
\item $\Delta E_\mathrm{c}=e\,U=\num{1.60e-19}\times\num{50e3}=\res{\qty{8,0e-15}{J}}$, soit \res{\qty{50}{keV}}.
\item Après $n$ intervalles, $E_\mathrm{c}=n\,e\,U$ et $v=\sqrt{\dfrac{2\,n\,e\,U}{m_\mathrm{p}}}$ : $v_1=\res{\qty{3,1e6}{m.s^{-1}}}$
et $v_{20}=\res{\qty{1,4e7}{m.s^{-1}}}$.
\item Aucune force notable : \res{\text{mouvement rectiligne uniforme}}.
\item $L=v\times\dfrac{T}{2}=\dfrac{v}{2f}$ : $L_1=\res{\qty{15}{cm}}$ et $L_{20}=\res{\qty{69}{cm}}$. La vitesse augmente à chaque
intervalle (comme $\sqrt{n}$), alors que la durée de traversée doit rester égale à $T/2$ : les tubes doivent s'allonger.
\end{questions}
\schema[corrige]{15-ex11}
\end{corrige}
\end{blocLivre}
\begin{blocLivre}{cor-15-12}
\begin{corrige}{12}{Le service au tennis}
\demarche{Tir horizontal ($\alpha=0$) depuis $h=\qty{2,80}{m}$ ; comparer la hauteur au filet et le point de chute.}
\begin{questions}
\item $x=45\,t$ et $z=-4{,}905\,t^2+2{,}80$ (origine au sol, sous le point de frappe).
\item Au filet ($x=\qty{11,9}{m}$, $t=\qty{0,264}{s}$) : $z=\qty{2,46}{m}$, elle passe largement. Au sol ($t=\qty{0,756}{s}$) :
$x=\res{\qty{34}{m}}$, bien au-delà du carré : \res{\text{service faute}}.
\item Pour tomber avant \qty{18,3}{m} en \qty{0,756}{s} : $v_0\leq\dfrac{18{,}3}{0{,}756}=\res{\qty{24}{m.s^{-1}}}$ (environ 87~km/h).
\item Avec $\alpha=-6{,}7^\circ$ : $v_x=\qty{49,7}{m.s^{-1}}$ et $v_z=\qty{-5,8}{m.s^{-1}}$. En $x=\qty{18,3}{m}$, $t=\qty{0,37}{s}$ et
$z\approx0$ : la balle tombe sur la ligne. Au filet, $z=\qty{1,13}{m}>\qty{0,914}{m}$ : \res{\text{elle franchit le filet}}.
\end{questions}
\end{corrige}
\end{blocLivre}
\begin{blocLivre}{cor-15-13}
\begin{corrige}{13}{L'imprimante à jet d'encre}
\demarche{Mouvement parabolique entre les plaques, puis rectiligne uniforme jusqu'au papier.}
\begin{questions}
\item $F=|q|\,E=\qty{1,5e-7}{N}$ et $P=m\,g=\qty{9,8e-10}{N}$ : $P/F=\num{0.0065}$, le poids est \res{\text{négligeable}}.
\item $a=\dfrac{|q|\,E}{m}=\res{\qty{1,5e3}{m.s^{-2}}}$ et $t=\dfrac{L}{v_0}=\res{\qty{5,0e-4}{s}}$.
\item $y_1=\frac12\,a\,t^2=\res{\qty{0,19}{mm}}$ et $v_y=a\,t=\res{\qty{0,75}{m.s^{-1}}}$.
\item Après les plaques : $\tan\theta=\dfrac{v_y}{v_0}=0{,}0375$, d'où une déviation supplémentaire $D\tan\theta=\qty{0,38}{mm}$ ;
au total \res{\qty{0,56}{mm}}. La déviation est proportionnelle à $q$ : l'imprimante \res{\text{règle la charge}} de chaque
goutte.
\end{questions}
\schema[corrige]{15-ex13}
\end{corrige}
\end{blocLivre}
\begin{blocLivre}{cor-15-14}
\begin{corrige}{14}{Le saut en longueur}
\demarche{Mouvement parabolique du centre de masse, de $z=\qty{1,10}{m}$ à $z=\qty{0,50}{m}$.}
\begin{questions}
\item $x=8{,}81\,t$ et $z=-4{,}905\,t^2+3{,}56\,t+1{,}10$ (en m et s).
\item $z_{\max}=1{,}10+\dfrac{3{,}56^2}{2\times9{,}81}=\res{\qty{1,75}{m}}$.
\item $z=0{,}50$ pour $t=\res{\qty{0,87}{s}}$, d'où $x=\res{\qty{7,6}{m}}$.
\item $\frac12\,m\,v^2=\frac12\,m\,v_0^2+m\,g\,(1{,}10-0{,}50)$, d'où $v=\res{\qty{10,1}{m.s^{-1}}}$.
\end{questions}
\end{corrige}
\end{blocLivre}
\begin{blocLivre}{cor-15-15}
\begin{corrige}{15}{La gouttelette de Millikan}
\demarche{Gouttelette immobile : la force électrique compense le poids, $|q|\,E=m\,g$ avec $E=U/d$.}
\begin{questions}
\item $m=\rho\times\frac43\pi r^3=\res{\qty{1,93e-15}{kg}}$ et $P=m\,g=\res{\qty{1,89e-14}{N}}$.
\item $\vec{E}$ vertical, de la plaque positive (en haut) vers la négative : \res{\text{vers le bas}}. La force électrique
doit être dirigée vers le haut, opposée à $\vec{E}$ : la charge est \res{\text{négative}}.
\item $|q|=\dfrac{m\,g\,d}{U}=\dfrac{\num{1.89e-14}\times0{,}0100}{394}=\res{\qty{4,81e-19}{C}}$, soit $3\,e$.
\item Toutes les charges sont des multiples entiers de $e$ ($2e$, $3e$, $5e$, $4e$) : la charge électrique est
\res{\text{quantifiée}}. Millikan a mesuré la \res{\text{charge élémentaire}} $e$.
\end{questions}
\end{corrige}
\end{blocLivre}
